\paragraph{One choice for both symbols at every nonzero address.}
For each $x\ne0$ let $P_{1,x},P_{2,x}$ be the two nonzero terminal
polynomials, of degrees $d_{1,x},d_{2,x}$, and let $H_{s,x}$ be their
highest homogeneous parts.  Work in the ordinary polynomial ring
in the terminal pair variables for this address.  Since it is a
domain, $H_{1,x}H_{2,x}\ne0$.  Choose a coefficient-one monomial
\[
 M_x=\prod_e X_e^{n_{x,e}},\qquad n_{x,e}\in\{0,1,2\},\qquad
 \sum_e n_{x,e}=d_{1,x}+d_{2,x}.
\]
For $n_{x,e}=2$ install one singular network $S_{x,e}$; it acts on
both symbol factors.  For $n_{x,e}=1$ install one varying network
$V^{\varepsilon_{x,e}}_{x,e}$.  Use disjoint auxiliary sets for all
these networks.  Let $\Gamma_y(\varepsilon)$ be the resulting graph
at address $y$, and consider the finite sum in $\mathbf F_2$
\[
 \sum_{\varepsilon}
       \prod_{y\ne0}\prod_{s=1}^2 A_s(\Gamma_y(\varepsilon)).
                                                               \tag{G.1a}
\]
The total degree available to all terminal polynomials is
$D=\sum_{y\ne0}(d_{1,y}+d_{2,y})$.

Each singular network at $x$ has zero constant at $x$, so its action
costs at least one difference in each of the two factors there, hence
at least two in (G.1a).  For a varying network, summing over its
$\varepsilon$ cancels the all-constant term.  Its first-degree
terms are independent of $\varepsilon$ away from its target, and
their difference at the target is $\partial_e$.  Thus every
surviving term costs at least one difference, and a cost-one term
acts at its target on exactly one of the two factors.  The sum of
all these minimum costs is exactly $D$.  Any higher-degree action,
or any additional action at a wrong address, exceeds $D$ and
annihilates at least one terminal polynomial.  Hence only the
minimum-cost terms survive.

At address $x$, a variable with exponent two is therefore
differentiated once in each factor; a variable with exponent one
is differentiated in exactly one factor, summed over the two
possibilities.  The resulting scalar is precisely
$[M_x](H_{1,x}H_{2,x})=1$.  Lower-degree parts cannot contribute.
Taking the product over all $x\ne0$ proves that (G.1a) equals one.
At least one choice of the finitely many $\varepsilon$ consequently
makes every factor in that product one.  This is one graph for both
symbols at every nonzero address.  Boolean reduction of
$H_{1,x}H_{2,x}$ before selecting $M_x$ would destroy this
exponent-two accounting and is not used.
Here is the stationary and one-symbol specialization, so that the
application in Proposition~\ref{ex:app-rough-companion} uses the same
proved construction. If all terminals are always active, attach
each column $h_i=Qe_i$ to any one fixed terminal. Since $h_i$ is
supported on its own group, their restricted sum is
$Q_J\mathbf1_J$ at every address. The two additional columns $c$
at a selected terminal pair cancel directly; no balancing terminal
of activation $\alpha+\beta$ is needed. The same singular program
therefore isolates that pair difference at its target address.
For one symbol, choose a coefficient-one top monomial of each
nonzero $P_x$ and put a separate singular program at $x$ for each
variable in that monomial. Its own programs already spend
$\deg P_x$ contraction degrees at $x$; any extra degree kills the
term. Thus exactly the selected top coefficient survives, and it
is one. A constant nonzero $P_x$ needs no program. This proves the
stationary one-symbol case as well as the moving one-symbol case.
In all cases the primary realization of
Lemma~\ref{ex:sec7-realization} changes only incidences with
the newly adjoined vertices. These are finite constructions for
each $b$, with no assertion of one infinite prime array.

\end{proof}

